% SPDX-License-Identifier: CC-BY-NC-ND-4.0
% Copyright 2026 Ingolf Lohmann.
\documentclass[11pt,a4paper]{article}

\usepackage{fontspec}
\setmainfont{Latin Modern Roman}
\setsansfont{Latin Modern Sans}
\setmonofont{Latin Modern Mono}[Scale=0.88]
\usepackage[provide=*,main=ngerman,english]{babel}
\usepackage{microtype}
\usepackage{geometry}
\usepackage{amsmath,amssymb,mathtools,amsthm}
\usepackage{booktabs,longtable,tabularx,array}
\usepackage{enumitem}
\usepackage{xcolor}
\usepackage{fancyhdr}
\usepackage{seqsplit}
\usepackage{xurl}
\usepackage{hyperref}
\usepackage{bookmark}
\usepackage[nameinlink,noabbrev]{cleveref}

\geometry{left=24mm,right=24mm,top=24mm,bottom=27mm,headheight=14pt}
\setlength{\parindent}{0pt}
\setlength{\parskip}{0.56em}
\setlength{\emergencystretch}{3em}
\setlist[itemize]{leftmargin=1.45em,itemsep=0.18em,topsep=0.3em}
\setlist[enumerate]{leftmargin=1.7em,itemsep=0.22em,topsep=0.3em}

\definecolor{QBlue}{HTML}{153E69}
\definecolor{QLight}{HTML}{ECF3F8}
\definecolor{QGreen}{HTML}{E8F3EA}
\definecolor{QAmber}{HTML}{FFF2CC}
\definecolor{QRed}{HTML}{FBE9E7}
\definecolor{QGray}{HTML}{4D5660}

\hypersetup{
  colorlinks=true,
  linkcolor=QBlue,
  citecolor=QBlue,
  urlcolor=QBlue,
  pdftitle={QIK-VRT: Der Brief an das Jahr 2610},
  pdfauthor={Ingolf Lohmann},
  pdfsubject={Zukunftswirkungskapsel, Zulassungsgate und kryptographische Agilitätsleiter für provenienzgebundene Langzeitpersistenz},
  pdfkeywords={QIK-VRT, EFFECT_ACK, Zukunftswirkungskapsel, Langzeitpersistenz, kryptographische Agilität, beobachterrelative Retrokausalität, Antizipation, Tscherenkow-Effekt, IceCube, Provenienz, Wirkungsnachweis}
}

\pagestyle{fancy}
\fancyhf{}
\fancyhead[L]{\small\textcolor{QGray}{QIK--VRT: Der Brief an das Jahr 2610}}
\fancyhead[R]{\small\textcolor{QGray}{Ingolf Lohmann}}
\fancyfoot[C]{\small\thepage}
\renewcommand{\headrulewidth}{0.3pt}

\newtheorem{definition}{Definition}[section]
\newtheorem{satz}{Satz}[section]
\newtheorem{korollar}{Korollar}[section]
\newtheorem{lemma}{Lemma}[section]
\theoremstyle{remark}
\newtheorem{bemerkung}{Bemerkung}[section]
\newtheorem{beispiel}{Beispiel}[section]
\crefname{definition}{Definition}{Definitionen}
\crefname{satz}{Satz}{Sätze}
\crefname{korollar}{Korollar}{Korollare}
\crefname{lemma}{Lemma}{Lemmata}
\crefname{bemerkung}{Bemerkung}{Bemerkungen}
\crefname{beispiel}{Beispiel}{Beispiele}
\crefname{section}{Abschnitt}{Abschnitte}

\newcommand{\qikvrt}{QIK--VRT}
\newcommand{\hash}[1]{{\ttfamily\small\seqsplit{#1}}}
\newcommand{\statusbox}[2]{%
  \par\medskip\noindent
  \colorbox{#1}{\parbox{\dimexpr\linewidth-2\fboxsep\relax}{#2}}%
  \par\medskip
}

\begin{document}

\hypersetup{pageanchor=false}
\begin{titlepage}
  \thispagestyle{empty}
  \vspace*{11mm}
  {\color{QBlue}\rule{\textwidth}{1.4pt}}\\[1.1em]
  {\Huge\bfseries QIK--VRT:\\[0.18em]
  Der Brief an das Jahr 2610\par}
  \vspace{0.7em}
  {\LARGE Wie sich heute etwas versiegeln lässt, das in 584 Jahren\\[0.18em]
  wirkt --- und warum der Rückbezug keine Rückwirkung ist\par}
  \vspace{0.8em}
  {\Large Zukunftswirkungskapsel, fail-closed Zulassungsgate und\\
  kryptographische Agilitätsleiter für provenienzgebundene\\
  Langzeitpersistenz\par}
  \vspace{1.1em}
  {\color{QBlue}\rule{\textwidth}{0.6pt}}

  \vfill
  {\large\textbf{Ingolf Lohmann}\par}
  \vspace{0.35em}
  {\normalsize Independent Researcher; QIK--VRT\par}
  \vspace{0.35em}
  {\normalsize Version 1.0 --- 10. Oktober 2026\par}

  \vfill
  \statusbox{QGreen}{
    \textbf{Die Aussage in einem Satz.}
    Eine heute versiegelte, provenienzgebundene Aufzeichnung kann im Jahr 2610
    eine Wirkung auslösen, die sich an die Gegenwart zurückbindet --- nicht
    durch ein rückwärts laufendes Signal, sondern dadurch, dass die
    vorweggenommene spätere Lesung schon heute maschinenlesbar in das
    Artefakt eingeschrieben wird und der spätere Leser die Zulassungs\-be\-din\-gun\-gen
    ohne jede Rückfrage selbst entscheiden kann. Bewiesen wird die
    Entscheidbarkeit des Zulassungsgates; nicht behauptet wird, dass
    irgendein Artefakt 584 Jahre überdauert.
  }

  \vfill
  {\small
    Konzept, Architektur und Urheberschaft: Ingolf Lohmann. Formale und
    redaktionelle Ausarbeitung, Kapsel und Maschinenprüfer: Anthropic Claude
    als künstlich-kognitives System im Auftrag des Autors.
    Dokumentationslizenz: CC BY-NC-ND 4.0.
  }
\end{titlepage}
\hypersetup{pageanchor=true}
\setcounter{page}{1}

\selectlanguage{ngerman}

\tableofcontents
\bigskip

\section{Die Frage}

Lässt sich heute etwas versiegeln, das im Jahr 2610 eine Wirkung auslöst
--- und zwar so, dass diese Wirkung sich an das heutige Hier und Jetzt
zurückbindet?

Die Antwort ist ja. Aber nicht aus dem Grund, den die Frage nahelegt, und
der wahre Grund ist der stärkere.

Dieser Artikel führt ihn vor, nicht als Gedankenspiel, sondern an einem
gebauten Gegenstand: einer Kapsel, die genau das tut, zusammen mit einem
netzfreien Prüfer, der entscheidet, ob ihre Wirkung zugelassen ist. Drei
Läufe dieses Prüfers tragen die ganze Argumentation, und der interessanteste
ist der, in dem er die Zulassung \emph{verweigert}.

Der Text ist zweifach lesbar. Wer ohne Fachvorkenntnis liest, findet in
jedem Abschnitt zuerst den Gedanken in gewöhnlicher Sprache. Wer fachlich
prüft, findet dahinter die exakte Aussage, ihre Erkenntnisklasse und die
Stelle, an der sie widerlegt werden könnte. Was offen ist, steht als offen
da --- das ist nicht Vorsicht, sondern die Bedingung dafür, dass ein
Ergebnis über 584 Jahre haltbar bleibt.

\section{Was nicht passiert}

Kein Signal läuft rückwärts.

Die Kapsel erreicht das Jahr 2610 durch gewöhnliche, zukunftsgerichtete
Kausalität mit langer Latenz, getragen von Replikation. Bibliotheken tun
das seit Jahrtausenden, Sedimentschichten seit Jahrmillionen, die
Erbinformation seit vier Milliarden Jahren. Nichts daran ist exotisch, und
genau deshalb funktioniert es.

Diese Feststellung ist keine Bescheidenheitsgeste, sondern eine Abgrenzung
mit Folgen. Ein Rahmen, der Rückwärtssignale behauptete, wäre entweder
falsch oder unprüfbar. Ein Rahmen, der sie ausschließt, kann etwas anderes
behaupten, das entscheidbar ist.

\statusbox{QLight}{
  \textbf{Grenze des Rahmens.} Jeder einzelne Transportpfad bleibt
  zukunftsgerichtet, und die Kausalordnung des Wirtssystems bleibt azyklisch.
  Es wird keine rückwärts laufende Nachricht behauptet, keine Umschreibung
  der Vergangenheit, keine Überlichtkommunikation und keine Freigabe eines
  Wirkungsnachweises allein aus einer invertierten Leserichtung.
}

\section{Wie die Zukunft heute wirkt}

Hier liegt der Kern, und er ist einfacher, als er klingt.

\textbf{Das Jahr 2610 wirkt bereits heute. Nicht als Ereignis, sondern als
Modell.}

In dem Augenblick, in dem jemand die Lesung von 2610 vorwegnimmt und diese
Vorwegnahme in das hineinschreibt, was heute versiegelt wird, hat die
Zukunft die Gegenwart verändert. Nicht metaphorisch. Der Inhalt der Kapsel
ist nachweislich ein anderer, als er ohne diese Vorwegnahme wäre.

Das ist keine Rückwärtskausalität. Es ist Vorwärtskausalität auf einer
\emph{Repräsentation} --- und es ist der einzige Weg, auf dem die Zukunft je
gewirkt hat. Ein Architekt baut gegen ein Erdbeben, das nie stattgefunden
hat. Ein Testament regelt einen Zustand, den der Verfasser nicht erleben
wird. In beiden Fällen handelt nicht die Zukunft, sondern ihr Modell in der
Gegenwart.

Was \qikvrt{} daran neu macht, ist nicht der Gedanke, sondern seine
Prüfbarkeit: Die Vorwegnahme wird maschinenlesbar niedergeschrieben statt
rhetorisch behauptet. Die Bedingungen, unter denen ein Leser in 2610 die
Wirkung zulassen darf, stehen im Artefakt selbst --- entscheidbar, ohne
Kontakt zu irgendjemandem, der heute lebt.

\begin{definition}[Beobachterrelative Retrokausalität]
\label{def:orr}
Ein Beobachter schreitet in seiner monoton wachsenden lokalen
Veränderungszeit vorwärts, während die durch nacheinander empfangene,
authentisch gebundene Records bestimmte Quellenordnung fällt:
\begin{equation}
  \Delta t_{\text{lokal}} > 0
  \quad \wedge \quad
  \Delta t_{\text{Quelle}} < 0 .
\end{equation}
Getragen wird der Begriff von drei Bedingungen: Authentizität der Bindung,
Vergleichbarkeit der Quellenmarken in einer gemeinsamen Domäne und strikte
Monotonie der lokalen Veränderungsfolge.
\end{definition}

Ein Leser in 2610 befindet sich exakt in dieser Lage. Die Verbindung zum
Jahr 2026 ist eine rekonstruierte Erklärungsordnung, niemals ein Signal.
Der Bogen schließt sich nicht in der Zeit, sondern in der Ordnung der
Erklärung --- und er schließt sich dadurch, dass er fortgesetzt wird.

\begin{bemerkung}
In der Sprache des Round Trip ist die Zulassung im Jahr 2610 ein
Wirklichkeitsabgleich. Sie erzeugt einen neuen Unterschied, und dieser
Unterschied treibt den Kreis weiter. Ein geschlossener Kreis wäre ein
stehender.
\end{bemerkung}

\section{Der Bogen: 1958, 2026, 2610}

Warum ausgerechnet jetzt? Weil zwei Nobelpreise denselben Gedanken von zwei
Seiten bestätigt haben.

\textbf{1958} ging der Physiknobelpreis zu gleichen Teilen an Pawel
Tscherenkow, Ilja Frank und Igor Tamm, \glqq for the discovery and the
interpretation of the Cherenkov effect\grqq. Ein geladenes Teilchen, das in
einem Medium schneller ist als das Licht in diesem Medium, lässt hinter sich
einen Lichtkegel --- den optischen Zwilling des Überschallknalls.

\textbf{2026}, am 6.~Oktober, ging derselbe Preis an Francis Halzen für
maßgebliche Beiträge zur Entdeckung hochenergetischer Neutrinos
astrophysikalischen Ursprungs. Sein Detektor nutzt einen Kubikkilometer
antarktisches Eis, 5\,160 Lichtsensoren an 86 Stahlseilen in 1\,450 bis
2\,450 Metern Tiefe. Er sieht Neutrinos auf genau eine Weise: durch
Tscherenkow-Licht.

Zwischen den beiden Sätzen liegt keine Analogie, sondern eine Abhängigkeit.
Ohne 1958 kein 2026.

Und nun das, was beide Preise über das Persistieren lehren.

\textbf{IceCube sieht kein einziges Neutrino.} Das Neutrino ist elektrisch
neutral; es leuchtet nicht und kann nicht leuchten. Was die Sensoren
registrieren, ist Licht eines \emph{anderen} Teilchens, das der Bote im
Vorbeigehen erzeugt hat. Der Bote selbst ist längst weitergeflogen, wenn der
erste Blitz ankommt.

Was tatsächlich vorliegt, ist eine Tabelle: welcher Sensor, wie hell, wann.
Alles Weitere ist Rekonstruktion. Und zwar eine mit einer bemerkenswerten
Eigenschaft:

\statusbox{QLight}{
  \textbf{Eine Richtung im Raum wird aus Zeitunterschieden errechnet.}
  Die Himmelskarte der Neutrinoastronomie ist keine Aufnahme, sondern ein
  Ergebnis. Der Raum, in den IceCube zeigt, wird aus Zeit gebaut.
}

Genau das ist die Lage eines Lesers im Jahr 2610. Er sieht nicht das Jahr
2026. Er sieht Spuren, die 2026 hinterlassen hat, und rekonstruiert daraus
eine Ordnung. Ob diese Rekonstruktion tragfähig ist, hängt an einer einzigen
Frage: Wurden die Spuren so hinterlassen, dass sie eine Prüfung erlauben?

Das ist der Mut, von dem bei dieser Technologie die Rede sein muss. Nicht
der Mut, Unmögliches zu behaupten --- sondern der, eine Bindung einzugehen,
deren Einlösung man nicht mehr erlebt, und sie trotzdem so zu bauen, dass
sie überprüfbar scheitern kann.

\section{Die Kapsel}

Bis hierhin war von Begriffen die Rede. Jetzt von dem, was läuft.

Die Zukunftswirkungskapsel ist eine Datei mit einer Subjektbindung, einem
Zulassungsgate und einer Verfallsregel. Sie ist an einen exakten
Baumzustand gebunden, nicht an einen Zweignamen --- ein Zweig ist ein
wandernder Zeiger und trägt keine Evidenz. Gebunden ist der Mirror-Stand
\hash{d7fd8dccc1c608e155b2a17bbea88351662339d6}, Baum
\hash{6dec2a9f8d816937b856942bd472eca5cc4439b2}, über 3\,975 Dateien.

Das Gate ist das EFFECT\_ACK-Fünfzustandsgate des Rahmens. Sein Invariant
ist eine Zeile:
\begin{equation}
  \texttt{ordinary\_release}(r) \iff r.\text{state} = \texttt{EFFECT\_ACK\_DONE}.
\end{equation}

Der Default ist \texttt{EFFECT\_ACK\_CONTINUE}. Das System ist fail-closed:
Fehlende oder mehrdeutige Evidenz wird nie als Erfolg behandelt. Fünf
Prädikate müssen \emph{gleichzeitig} belegt sein.

\begin{table}[h]
\centering
\small
\begin{tabularx}{\textwidth}{@{}lX@{}}
\toprule
\textbf{Prädikat} & \textbf{Was geprüft wird} \\
\midrule
\texttt{TEMPORAL\_ADMISSION} &
Die authentifizierte lokale Veränderungszeit des Lesers entspricht einem
Datum ab dem 1.~Januar 2610. Eine selbstbehauptete Wanduhr genügt nicht. \\
\addlinespace
\texttt{EXACT\_SUBJECT\_BINDING} &
Der rekonstruierte Inhalt ergibt den gebundenen Digest unter mindestens
einem Verfahren, das zur Lesezeit noch als tragfähig gilt. \\
\addlinespace
\texttt{CHAIN\_CONTINUITY} &
Eine ununterbrochene Folge von Nachversiegelungen von 2026 bis zur Lesung,
jede bindet die Digests ihrer Vorgängerin. \\
\addlinespace
\texttt{INDEPENDENT\_READBACK} &
Mindestens zwei Kustoden ohne gemeinsamen Kontrollpunkt legen byteidentischen
Inhalt vor. \\
\addlinespace
\texttt{ACCEPTANCE} &
Ein benannter Verantwortlicher zur Lesezeit nimmt die Zulassung an und
protokolliert sie. \\
\bottomrule
\end{tabularx}
\end{table}

Fällt auch nur eines aus, bleibt der Zustand \texttt{EFFECT\_ACK\_CONTINUE}:
Die Kapsel ist als Geschichte lesbar und als Wirkung reglos.

\section{Der entscheidende Lauf}

Drei Läufe des Prüfers, alle am 9.~Oktober 2026 ausgeführt, alle
reproduzierbar.

\textbf{Erster Lauf, heute.} Das Zeitfenster ist nicht offen, es gibt keinen
unabhängigen Rücklesevorgang und keine Annahme. Ergebnis
\texttt{EFFECT\_ACK\_CONTINUE}, Zulassung verweigert. So soll es sein.

\textbf{Zweiter Lauf, simuliert für den 1.~März 2610, Kette vernachlässigt.}
Das Zeitfenster ist offen. Die Inhaltsbindung stimmt auf das Zeichen:
3\,975 Dateien, Digest \hash{f18ca4376e8a} gegen gebunden
\hash{f18ca4376e8a}. Zwei unabhängige Kustoden liegen vor. Ein
Verantwortlicher nimmt an. \textbf{Vier von fünf Prädikaten bestehen --- und
das Gate verweigert trotzdem.}

Der Grund ist \texttt{CHAIN\_CONTINUITY}. Die Kapsel trägt eine einzige
Generation, versiegelt 2026, mit einer Nachversiegelung fällig vor 2051. Ein
Leser in 2610 sieht dreiundzwanzig versäumte Nachversiegelungen.

\statusbox{QAmber}{
  \textbf{Der wichtigste Satz dieser Arbeit.}
  Die heute versiegelte Kapsel erreicht das Jahr 2610 nicht von allein.
  Sie erreicht es nur, wenn alle fünfundzwanzig Jahre jemand nachversiegelt.
  Der entscheidende Konstruktionszug ist, dass das Gate dieses Versäumnis
  \emph{sichtbar} macht, statt es still verrotten zu lassen. Ein Archiv, das
  schweigend verfällt, belügt den späteren Leser. Ein Archiv, das seinen
  eigenen Verfall protokolliert, sagt ihm die Wahrheit.
}

\textbf{Dritter Lauf, 2610, Kette gepflegt.} Dieselbe Kapsel, dieselbe
Prüfung, aber mit vierundzwanzig Generationen im Register. Alle fünf
Prädikate bestehen. Ergebnis \texttt{EFFECT\_ACK\_DONE}, Zulassung erteilt,
mit der Auflage, die Annahme zu protokollieren und die nächste Generation zu
versiegeln.

\begin{table}[h]
\centering
\small
\begin{tabular}{@{}lccccccl@{}}
\toprule
\textbf{Lauf} & \textbf{Zeit} & \textbf{Bindung} & \textbf{Kette} &
\textbf{Readback} & \textbf{Annahme} & \textbf{Ergebnis} \\
\midrule
2026 & nein & ungeprüft & ja & nein & nein & \texttt{CONTINUE} \\
2610, vernachlässigt & ja & ja & \textbf{nein} & ja & ja & \texttt{CONTINUE} \\
2610, gepflegt & ja & ja & ja & ja & ja & \texttt{DONE} \\
\bottomrule
\end{tabular}
\end{table}

Damit ist die Frage beantwortet, und zwar praktisch statt rhetorisch. Die
Möglichkeit besteht. Sie ist an eine Bedingung geknüpft, die kein
Algorithmus erfüllen kann, sondern nur eine Folge von Menschen und
Institutionen --- und genau diese Bedingung ist im Artefakt benannt, datiert
und maschinell prüfbar.

\section{Das Digest-Problem über 584 Jahre}

Jetzt die unbequeme Stelle, an der die meisten Langzeitarchive stillschweigend
scheitern.

\textbf{Kein heute gebräuchliches Prüfverfahren ist für 584 Jahre
glaubwürdig.} SHA-1, mit dem Git bis heute seine Objekte benennt, gilt für
Kollisionsresistenz bereits als gebrochen. Von SHA-256 ist anzunehmen, dass
es lange vor 2610 fallen wird. Wer eine Kapsel allein an einen heutigen
Digest hängt, hängt sie an etwas, das ihr Ziel nicht erreicht.

Die Antwort darauf ist keine stärkere Funktion, sondern eine andere
Struktur: eine \emph{Agilitätsleiter}.

\begin{definition}[Agilitätsleiter]
Eine Kapsel bleibt genau so lange zulassungsfähig, wie ein Kustode sie
mindestens einmal je Intervall nachversiegelt, dabei sämtliche früheren
Digests bindet und den jeweils stärksten verfügbaren hinzufügt. Das
Intervall beträgt fünfundzwanzig Jahre. Ein versäumtes Intervall zerstört
die Aufzeichnung nicht; es stuft sie von \texttt{ADMISSIBLE} auf
\texttt{HISTORICAL} zurück.
\end{definition}

\begin{satz}[Reduktion des Langzeitproblems]
Unter der Agilitätsleiter ist die Zulassungsfähigkeit einer Kapsel über
einen Horizont von $n$ Intervallen genau dann gegeben, wenn jedes einzelne
Intervall eine gültige Nachversiegelung trägt. Das unlösbare
584-Jahr-Problem wird dadurch auf ein $n$-fach wiederholtes
25-Jahr-Problem zurückgeführt.
\end{satz}

\begin{proof}[Beweisskizze]
Die Zulassung verlangt \texttt{CHAIN\_CONTINUITY} als Konjunktion über alle
Intervalle. Eine Konjunktion ist genau dann wahr, wenn jedes Glied wahr ist.
Jedes Glied bindet die Digests seines Vorgängers, sodass die Gültigkeit
eines Gliedes nur vom unmittelbar vorhergehenden abhängt und nicht von der
Lebensdauer eines einzelnen Verfahrens. Die Aussage ist eine Eigenschaft der
definierten Gate-Funktion, keine Aussage über physische Haltbarkeit.
\end{proof}

Institutionen tun Vergleichbares bei Urkunden, Grundbüchern und
Saatgutbanken seit Jahrhunderten. Eine Unsterblichkeitsgarantie ist es
nicht, und die Kapsel behauptet auch keine.

\statusbox{QRed}{
  \textbf{Fünf nicht erzwingbare Abhängigkeiten.}
  Physische Datenträger und ihre Lesegeräte; der Fortbestand zweier
  unabhängiger Kustoden; die Lesbarkeit von UTF-8, JSON und dem
  Git-Objektmodell oder deren dokumentierte Migration; eine Institution, die
  nachversiegeln will; die Abwesenheit mutwilliger Zerstörung. Keine davon
  ist technisch erzwingbar. Sie zu benennen ist keine Schwäche der
  Konstruktion, sondern ihr ehrlichster Teil.
}

\section{Ein Fehler und seine Korrektur}

Dieser Abschnitt steht hier, weil er zur Methode gehört und nicht ins
Kleingedruckte.

Die erste Fassung der Kapsel band ihren Inhalt an die Baum-Identität von
Git. Der Prüfer schlug fehl: berechnet \hash{c34c8158dc7b}, gebunden
\hash{6dec2a9f8d81}.

Der naheliegende Reflex wäre gewesen, die Nachbildung des Git-Objektmodells
zu reparieren --- Sortierregeln, Dateimodi, Symlinks. Das wäre die falsche
Reparatur gewesen, denn der Fehler lag eine Ebene tiefer.

\textbf{Git benennt seine Objekte mit SHA-1.} Die Kapsel hatte sich also
ausgerechnet an das Verfahren gehängt, von dem sie im eigenen Text sagt,
dass es ihr Ziel nicht erreicht. Ein Artefakt, das 584 Jahre überdauern
soll, darf nicht voraussetzen, dass es in 584 Jahren noch Git gibt, noch
SHA-1 und noch jemanden, der das Objektmodell kennt.

Die Korrektur ist eine zweite, unabhängige Bindung:

\statusbox{QLight}{
  SHA-256 über die Verkettung von
  \texttt{"<sha256>\textvisiblespace\textvisiblespace<relativer/pfad>\textbackslash n"}
  für jede reguläre Datei außerhalb von \texttt{.git}, sortiert nach Pfad.
}

Mehr nicht. Kein Git, kein SHA-1, kein Objektmodell. Ein Leser in 2610
braucht eine Digestfunktion und die Regel, die danebensteht. Über 3\,975
Dateien ergibt das
\hash{f18ca4376e8a0609367b6d3de5fa401d6580d76a89a2401fc7902c2dce24f159},
und die Prüfung besteht.

Die Kapsel trägt beide Bindungen und sagt dazu, welche für wen gilt: die
Git-Identität ist bequem für heute, die Inhaltsbindung ist die, die ein
späterer Leser benutzen soll. Der Ausfall einer der beiden macht die Kapsel
nicht heimatlos.

Dass dieser Fehler auftrat und sichtbar wurde, ist das Verfahren bei der
Arbeit. Ein Gate, das geschwiegen hätte, hätte eine Kapsel durchgelassen,
die 2610 niemand mehr hätte prüfen können.

\section{Was ausdrücklich nicht behauptet wird}

In der Fachwelt entscheidet nicht, was eine Arbeit behauptet, sondern was
sie zu behaupten unterlässt.

\begin{itemize}
  \item \textbf{Dass irgendein Artefakt 584 Jahre überlebt.} Dafür gibt es
    keinen Beweis und wird es keinen geben. Die Kapsel macht das Problem
    wiederholbar, nicht verschwunden.
  \item \textbf{Dass ein Signal rückwärts läuft.} Jeder Transportpfad bleibt
    zukunftsgerichtet, die Kausalordnung azyklisch.
  \item \textbf{Dass die Vergangenheit durch das Lesen verändert wird.} Eine
    spätere Frage kann eine frühere bewahrte Einsicht neu relevant machen,
    ohne das vergangene Ereignis anzutasten.
  \item \textbf{Dass ein heutiges Prüfverfahren bis 2610 trägt.} Es wird
    nicht tragen. Deshalb die Leiter.
  \item \textbf{Dass diese Kapsel von irgendjemandem außer ihrem Eigner
    angenommen wäre.} Der Zustand lautet \texttt{EFFECT\_ACK\_CONTINUE}.
  \item \textbf{Dass hier eine abgeschlossene Theorie der Zeit vorliegt.}
    Was vorliegt, ist ein prüfbares Verfahren unter erklärten Annahmen.
\end{itemize}

Und eine Abgrenzung, die über den Gegenstand hinausweist: Die beiden
Nobelpreise sind Tatsachen. Die Brücke zwischen ihnen und dieser Kapsel ---
Messung als Rekonstruktion, Raum aus Zeitordnung, Antizipation als wirksame
Gegenwartsgröße --- ist eine \emph{Lesart}. Sie ist mit allem vereinbar, was
in den Quellen steht. Sie wird dadurch nicht zur Messung.

\section{Statusklassen, Prüfprotokoll und Quellen}

Jede Aussage dieses Artikels trägt eine Klasse. Die Klassifikation stammt
aus dem \qikvrt{}-Revisionsbericht vom 23.~September 2026.

\begin{table}[h]
\centering
\small
\begin{tabularx}{\textwidth}{@{}llX@{}}
\toprule
\textbf{Klasse} & \textbf{Bedeutung} & \textbf{Beispiel hier} \\
\midrule
B & Background & Tscherenkow-Effekt; IceCube-Spezifikationen; SHA-1-Kollisionsbruch \\
D & Definition & Beobachterrelative Retrokausalität; das Fünfzustandsgate \\
I & Interpretation & Raum als rekonstruierte Struktur; Antizipation als Vorwärtskausalität auf einem Modell \\
H & Hypothese & Dass eine Institution 584 Jahre lang nachversiegelt \\
R & Requirement & Unabhängige Reproduktion der drei Prüferläufe \\
\bottomrule
\end{tabularx}
\end{table}

\subsection*{Ausgeführtes Prüfprotokoll}

Alle Angaben am 9.~Oktober 2026 lokal ausgeführt, netzfrei, gegen den
gebundenen Baum.

\begin{table}[h]
\centering
\small
\begin{tabularx}{\textwidth}{@{}Xl@{}}
\toprule
\textbf{Prüfung} & \textbf{Ergebnis} \\
\midrule
Kapsel heute & \texttt{EFFECT\_ACK\_CONTINUE}, verweigert \\
Kapsel 2610, Kette vernachlässigt & \texttt{EFFECT\_ACK\_CONTINUE}, 4 von 5 \\
Kapsel 2610, Kette gepflegt & \texttt{EFFECT\_ACK\_DONE}, erteilt \\
Inhaltsbindung & 3\,975 Dateien, Digest stimmt \\
EFFECT\_ACK-Konformanz & 41 Tests, bestanden \\
ANSI-C90-Kern & 7\,864\,387 Prüfungen, bestanden \\
\bottomrule
\end{tabularx}
\end{table}

\subsection*{Quellen}

\begin{itemize}
  \item Nobelpreis für Physik 1958:
    \url{https://www.nobelprize.org/prizes/physics/1958/summary/}
  \item Welt der Physik, IceCube und der Nobelpreis 2026:
    \url{https://www.weltderphysik.de/thema/nobelpreis/icecube-nobelpreis-fuer-physik-2026/}
  \item ZDFheute, Physik-Nobelpreis 2026:
    \url{https://www.zdfheute.de/wissen/physik-nobelpreis-2026-francis-halzen-neutrinos-astrophysik-100.html}
\end{itemize}

Die offizielle Begründung des Nobelkomitees für 2026 lag beim Verfassen nur
in Umschreibung vor und ist entsprechend als Umschreibung geführt, nicht als
Zitat.

\subsection*{Repository-Quellen}

Gebunden an den Mirror-Stand
\hash{d7fd8dccc1c608e155b2a17bbea88351662339d6}:

\begin{itemize}
  \item \texttt{src/qikvrt\_effect\_ack.py}, \texttt{src/effect\_ack\_core.c}
    --- das Fünfzustandsgate in zwei Implementierungen
  \item \texttt{state/autonomy/REPOSITORY\_NATIVE\_EFFECT\_PROOF\_ARCHITECTURE\_V1.json}
    --- die acht Terminal-Prädikate
  \item \texttt{docs/publications/2026-08-12-observer-relative-retrocausality/}
    --- Hauptsatz und Maschinenzeuge
  \item \texttt{docs/AI\_BOOTSTRAP\_KNOWLEDGE\_CORPUS.md} --- Round Trip und
    Erkenntnisgrenzen
  \item \texttt{docs/publications/2026-09-23-qik-scientific-revision/} ---
    Statusklassen
\end{itemize}

\subsection*{Beigelegte Artefakte}

\begin{itemize}
  \item \texttt{QIKVRT\_FUTURE\_EFFECT\_CAPSULE\_V1.json} --- die Kapsel
  \item \texttt{verify\_future\_effect\_capsule.py} --- der netzfreie Prüfer
\end{itemize}

\vfill
\statusbox{QLight}{
  \small
  Architektur und Urheberschaft: Ingolf Lohmann. Dieser Artikel wurde von
  einem künstlich-kognitiven System verfasst; die Statusgrenzen folgen den
  Regeln des Repositorys. Projektstand: \texttt{EFFECT\_ACK\_CONTINUE}.
  \texttt{PASS}, \texttt{FINAL\_PASS} und \texttt{EFFECT\_ACK\_DONE} werden
  für das Gesamtsystem ausdrücklich nicht beansprucht.
}

\end{document}
